% !TeX program = lualatex
%==============================================================================
%  Leant — a Djex-based synthesis REPL for Lean 4
%  An overview and tutorial, with a detailed exposition of :synth.
%
%  Leant is an experimental project in active development; this document
%  describes the state of the tool at the date on the title page.
%
%  Build:  lualatex Leant.tex   (twice, for the table of contents)
%
%  The visual design (palette, fonts, code boxes) follows the typeset
%  edition of "Counterexamples in Type Systems".
%==============================================================================
\documentclass[11pt,oneside]{article}

\usepackage[letterpaper,
            textwidth=5.65in, textheight=8.5in,
            centering, headsep=16pt, footskip=28pt]{geometry}

\usepackage{amsmath}
\usepackage{fontspec}
\usepackage{unicode-math}
\usepackage{microtype}
\usepackage{xcolor}
\usepackage{fvextra}
\usepackage[most]{tcolorbox}
\usepackage{enumitem}
\usepackage{titlesec}
\usepackage{fancyhdr}
\usepackage{needspace}
\usepackage[hang,flushmargin,bottom]{footmisc}
\usepackage{xurl}
\usepackage{hyperref}
\usepackage{bookmark}

%------------------------------------------------------------------ fonts ----
\setmainfont{LibertinusSerif}[
  Extension    = .otf,
  UprightFont  = *-Regular, ItalicFont = *-Italic,
  BoldFont     = *-Bold,    BoldItalicFont = *-BoldItalic,
  Ligatures    = TeX ]
\setsansfont{LibertinusSans}[
  Extension    = .otf,
  UprightFont  = *-Regular, ItalicFont = *-Italic, BoldFont = *-Bold,
  Ligatures    = TeX ]
\setmonofont{DejaVuSansMono}[
  Extension    = .ttf,
  UprightFont  = *, ItalicFont = *,
  ItalicFeatures = {FakeSlant = 0.2},
  BoldFont     = *-Bold, BoldItalicFont = *-Bold,
  BoldItalicFeatures = {FakeSlant = 0.2},
  Scale        = 0.83 ]
\setmathfont{LibertinusMath-Regular.otf}

%----------------------------------------------------------------- colours ---
%  A calm, paper-like palette: sage greens, warm off-whites, near-black.
\definecolor{ink}      {HTML}{0B0B0B}  % body text
\definecolor{moss}     {HTML}{3E5C43}  % links, rules, section numbers
\definecolor{sage}     {HTML}{D7E8CD}  % code label ground
\definecolor{sagepale} {HTML}{F8FBF1}  % code ground
\definecolor{sageline} {HTML}{DFE4D7}  % hairlines
\definecolor{parchment}{HTML}{F4F3EB}  % warm ground
\definecolor{stone}    {HTML}{ADADAC}  % dimmed text
\definecolor{lichen}   {HTML}{77836A}  % tag line
\definecolor{cmtcol}   {HTML}{7A8A76}  % code comments
\definecolor{kwcol}    {HTML}{33523B}  % code keywords
\definecolor{strcol}   {HTML}{8A5A44}  % code strings
\color{ink}

%---------------------------------------------------------------- headings ---
\titleformat{\section}
  {\normalfont\LARGE\bfseries\color{ink}}
  {\color{moss}\thesection}{0.7em}{}
  [\vspace{2pt}{\color{sageline}\titlerule[0.8pt]}]
\titlespacing*{\section}{0pt}{0pt}{1.1\baselineskip}

\titleformat{\subsection}{\normalfont\large\bfseries\color{ink}}
  {\color{moss}\thesubsection}{0.6em}{}
\titlespacing*{\subsection}{0pt}{1.4\baselineskip}{0.45\baselineskip}

\setcounter{secnumdepth}{2}
\setcounter{tocdepth}{2}

%----------------------------------------------------------------- running ---
\pagestyle{fancy}
\fancyhf{}
\renewcommand{\headrulewidth}{0pt}
\fancyhead[L]{\footnotesize\itshape\color{stone}\nouppercase{\leftmark}}
\fancyhead[R]{\footnotesize\color{stone}\thepage}
\fancypagestyle{plain}{\fancyhf{}%
  \fancyhead[R]{\footnotesize\color{stone}\thepage}}
\renewcommand{\sectionmark}[1]{\markboth{#1}{}}

%------------------------------------------------------------------- lists ---
\setlist[itemize]{leftmargin=1.35em, itemsep=0.45\baselineskip,
                  parsep=0.35\baselineskip, topsep=0.5\baselineskip,
                  label={\color{moss}\rule[0.28ex]{0.36em}{0.36em}}}

%------------------------------------------------------------------- misc ----
\newcommand{\code}[1]{\texttt{#1}}
\newcommand{\cmd}[1]{\texttt{\color{kwcol}#1}}

\widowpenalty=10000
\clubpenalty=10000
\raggedbottom
\setlength{\emergencystretch}{2em}

%---------------------------------------------------------------- code box ---
%  REPL transcripts are typeset by fvextra rather than listings: the
%  transcripts are full of non-ASCII Lean syntax, and listings under
%  LuaTeX reorders a multibyte character in front of whatever precedes
%  it.  The session environment writes its body verbatim to a scratch
%  file and reads it back inside the box.
\tcbset{cebox/.style={
  enhanced, breakable,
  colback=sagepale, colframe=sageline, boxrule=0.5pt, arc=2pt,
  left=7pt, right=7pt, top=5pt, bottom=5pt,
  toptitle=1pt, bottomtitle=1pt,
  coltitle=ink, colbacktitle=sage,
  fonttitle=\sffamily\scriptsize\bfseries,
  attach boxed title to top left={xshift=9pt, yshift=-\tcboxedtitleheight/2},
  boxed title style={boxrule=0pt, arc=1.5pt, left=5pt, right=5pt,
                     top=1.5pt, bottom=1.5pt},
  before skip=\medskipamount, after skip=\medskipamount }}

\newtcolorbox{sessionouter}{cebox, title={session}}

\newenvironment{session}
 {\VerbatimEnvironment\begin{VerbatimOut}{\jobname.session.tmp}}
 {\end{VerbatimOut}%
  \begin{sessionouter}%
  \VerbatimInput[fontsize=\footnotesize, breaklines=true,
    breaksymbolleft={\textcolor{stone}{\footnotesize$\hookrightarrow$}},
    formatcom=\color{ink}, tabsize=4]{\jobname.session.tmp}%
  \end{sessionouter}}

% the "experimental" notice
\newtcolorbox{notice}[1]{enhanced, breakable,
  colback=parchment, colframe=sageline, boxrule=0.5pt, arc=2pt,
  left=8pt, right=8pt, top=6pt, bottom=6pt,
  coltitle=ink, colbacktitle=sage,
  fonttitle=\sffamily\scriptsize\bfseries, title={#1},
  attach boxed title to top left={xshift=9pt, yshift=-\tcboxedtitleheight/2},
  boxed title style={boxrule=0pt, arc=1.5pt, left=5pt, right=5pt,
                     top=1.5pt, bottom=1.5pt},
  before skip=\medskipamount, after skip=\medskipamount }

%------------------------------------------------------------ command table --
\newenvironment{cmdtable}
  {\par\begingroup\small
   \setlength{\parindent}{0pt}\setlength{\parskip}{2.5pt}}
  {\endgroup\par\vspace{0.3\baselineskip}}
\newcommand{\cmdrow}[2]{%
  \makebox[11.5em][l]{\code{#1}}\parbox[t]{\dimexpr\linewidth-12em}{#2}\par}

%--------------------------------------------------------------- hyperlinks --
\hypersetup{
  colorlinks = true,
  linkcolor  = moss,
  urlcolor   = moss,
  citecolor  = moss,
  pdftitle   = {Leant — a Djex-based synthesis REPL for Lean 4},
  pdfauthor  = {The Leant project},
  pdfsubject = {Verified term synthesis for Lean 4, in an interactive REPL},
  pdfkeywords= {Lean 4, REPL, term synthesis, Djinn, LJT, Exference},
  bookmarksnumbered = true }
\urlstyle{same}

\begin{document}
%------------------------------------------------------------------ title ---
\begin{titlepage}
\thispagestyle{empty}
\vspace*{\stretch{2}}
\begin{center}
  {\color{sageline}\rule{0.55\textwidth}{0.8pt}}\\[1.7em]
  {\fontsize{34}{40}\selectfont Leant}\\[0.55em]
  {\Large a Djex-based synthesis REPL for Lean 4}\\[1.5em]
  {\color{sageline}\rule{0.55\textwidth}{0.8pt}}\\[2.2em]
  {\large \emph{an overview and tutorial, with a detailed tour of \code{:synth}}}
\end{center}
\vspace{\stretch{3}}
\begin{center}
  \small\color{lichen}
  An \textbf{experimental} tool in active development.
  Interfaces, output, and capabilities change frequently.\\[0.6em]
  \href{https://github.com/VladimirReshetnikov/Leant}{github.com/VladimirReshetnikov/Leant}\quad
  August 2026.
\end{center}
\vspace*{\stretch{1}}
\end{titlepage}

\clearpage
{\hypersetup{linkcolor=ink}%
 \microtypesetup{protrusion=false}%
 \tableofcontents
 \microtypesetup{protrusion=true}}
\clearpage

%==============================================================================
\section{Introduction}\label{sec:intro}

Leant is an interactive read--eval--print loop for Lean~4, in the
spirit of GHCi: you type expressions and they are evaluated, you type
declarations and they enter the session, and a family of colon-commands
gives you type queries, documentation, search, interactive proving,
and automatic \emph{term synthesis} with machine-checked answers.
The synthesis command takes up half of this document.

Lean~4 itself is normally driven from an editor: the language server
shows goals and diagnostics as you edit a file. Leant complements that
workflow with a conversational one. It is aimed at exploration (trying a
lemma, poking at a definition, asking \emph{is this type even
inhabited?}), where the unit of work is a line, not a file.

\begin{notice}{experimental}
Leant is an \textbf{experimental project in active development}. It is
a research vehicle, not a supported product: commands appear, change
shape, and disappear between commits; output formats are not stable;
and some features are bounded approximations by design. Every
transcript in this document was produced by the tool at the date on
the title page and may not match the tool you build tomorrow. When
this document and \code{:help} disagree, believe \code{:help}.
\end{notice}

Leant is implemented in Haskell and speaks the backend protocol
directly; term synthesis (\S\ref{sec:synth}) embeds the Djex synthesis
library in-process.

Under the hood, Leant drives the community
\code{leanprover-community/repl} binary over a JSON protocol: every
input becomes a command against a persistent Lean environment, and
every claim Leant makes, including every synthesized term, is
checked by that backend before it is shown. Sessions survive backend
crashes (the session history replays automatically), can be saved and
restored (\code{:pickle}/\code{:unpickle}, including synthesis-ready
companions for Leant-created snapshots), and can be recorded as
transcripts.

%==============================================================================
\section{Getting started}\label{sec:start}

Build the binary once with \code{cabal build exe:leant} (GHC 9.12.4;
the vendored \code{lib/Djex} submodule must be initialized), then run
\code{leant.cmd}, or the built executable directly. Leant finds the
Lean backend the same way LeanInteract's cache laid it out, or accepts
an explicit \code{-{}-repl-exe}/\code{LEANT\_BACKEND}. Without a Lake
project it runs against the plain Lean~4 standard library, which is
also the least memorable thing about it: no imports, subsecond
startup.

A session begins as a calculator and escalates quietly:

\begin{session}
λ> 2 + 2
4
λ> it * 10
40
λ> def double (n : Nat) : Nat := n + n
λ> double 21
42
λ> :type double
double : Nat → Nat
\end{session}

Expressions are tried with \code{\#eval} and fall back to
\code{\#check}; the last evaluated value is bound as \code{it},
GHCi-style. Declarations (\code{def}, \code{theorem},
\code{inductive}, \dots) extend the session environment, and any
\code{\#}-command (\code{\#eval}, \code{\#check}, \code{\#print
axioms}) passes straight through. Multi-line input opens automatically
when a line is syntactically incomplete (type a \code{structure}
header and just keep going), and an empty line submits it.

%==============================================================================
\section{The commands}\label{sec:commands}

\subsection{Session and evaluation}
\begin{cmdtable}
\cmdrow{:help, :h, :?}{show the command summary}
\cmdrow{:quit, :q}{exit}
\cmdrow{:type EXPR, :t}{show the type of \code{EXPR} (\code{\#check})}
\cmdrow{:info NAME, :i}{show the definition of \code{NAME} (\code{\#print})}
\cmdrow{:load FILE, :l}{reset the session and load a \code{.lean} file}
\cmdrow{:reload, :r}{reload the last loaded file}
\cmdrow{:import MOD}{add an import (rebuilds the session; unavailable over an opaque snapshot)}
\cmdrow{:imports}{list active imports, or imports restored by \code{:reset}}
\cmdrow{:undo}{revert the last state-changing command}
\cmdrow{:reset}{clear definitions or a snapshot base (keeps configured imports)}
\cmdrow{:history}{show commands after the current import/snapshot base}
\cmdrow{:set OPT VAL}{a Lean \code{set\_option} that persists in the session}
\end{cmdtable}

\subsection{Discovery}
\begin{cmdtable}
\cmdrow{:browse [NS]}{list declarations in a namespace, or the session's own}
\cmdrow{:browse!\ NS}{\dots including compiler-generated auxiliaries}
\cmdrow{:doc NAME}{show a docstring}
\cmdrow{:search TEXT}{search declaration names, case-insensitively}
\cmdrow{:search?\ TYPE}{proof search: what proves \code{TYPE}? (via \code{exact?})}
\end{cmdtable}

\subsection{Proving and synthesis}
\begin{cmdtable}
\cmdrow{:prove [PROP]}{prove \code{PROP} interactively, tactic by tactic
  (without an argument: resume the last \code{sorry})}
\cmdrow{:synth TYPE}{synthesize verified terms of \code{TYPE};
  candidates are bound as \code{it1} (= \code{it}), \code{it2}, \dots
  (\S\ref{sec:synth})}
\cmdrow{:set synth-engine E}{\code{djinn} (default) $\mid$
  \code{exference} $\mid$ \code{both}}
\cmdrow{:set synth-steps N}{Exference's step budget (default 4096)}
\cmdrow{:set synth-classical B}{offer classical candidates for
  constructively refuted goals (\code{on}$\mid$\code{off}, default on)}
\cmdrow{:set synth-library B}{library premises for recursive
  inductives (\code{on}$\mid$\code{off}, default on;
  \S\ref{sec:synth-library})}
\end{cmdtable}

\subsection{Persistence and utility}
\begin{cmdtable}
\cmdrow{:transcript [F|on|off]}{record a transcript of the session}
\cmdrow{:timestamps [on|off]}{timestamp transcript entries}
\cmdrow{:pickle FILE}{save the environment and synthesis companion}
\cmdrow{:unpickle FILE}{restore an environment as a new undo/history base}
\cmdrow{:env}{show the backend environment id}
\cmdrow{:time}{toggle per-command timing}
\cmdrow{:!\ CMD}{run a shell command}
\end{cmdtable}

%==============================================================================
\section{Interactive proving}\label{sec:prove}

\code{:prove} turns the prompt into a tactic loop against the
backend's proof-state protocol. Every line is a tactic; \code{:undo}
pops one step; \code{:qed NAME} assembles the accumulated script into
a theorem and replays it into the session, so the result is usable
afterwards.  Each displayed proof state is followed by an automatically
discovered tactic suggestion that Lean has already checked.  The
candidate probes are shaped by the goal and its hypotheses: an
\code{intro} suggestion names the binders it would introduce, a
disjunction goal is probed with \code{left}/\code{right}, a
hypothesis whose type a single step can take apart is probed with
\code{cases} or \code{obtain}, and a data-typed variable the goal
mentions is probed with \code{induction}.  The search prefers a
candidate that closes the goal outright and annotates the suggestion
accordingly (\code{closes the goal}, \code{splits into 2 goals}); when
no single tactic closes it, a second phase chains quick finishers onto
the best progressing candidates (\code{constructor <;> omega}, or an
\code{obtain} step chained with the \code{exact} it enables), so even
the opening suggestion is often a complete checked proof.  The chains
also try \code{simp\_all}, unfolding the definitions the goal mentions
and calling in \code{omega} on whatever arithmetic remains, which is
how an induction suggestion can arrive as a finished proof:
\code{induction l <;> simp\_all [myLen]}, or
\code{induction n <;> simp\_all [double] <;> omega}.
The suggestion is advisory---it does not advance the state or enter the
script---and \code{:suggest} reprints it without rerunning the search.
Below, the very first suggestion is already a complete proof; the session
ignores it and proceeds step by step instead.
Proof-state identifiers are local to one backend process; if it stops,
Leant leaves prove mode and prints the accumulated script rather than
submitting a stale identifier after session reconstruction.

\begin{session}
λ> :prove ∀ p q : Prop, (p → q) → p → q ∧ p
entering prove mode — type tactics; :help for commands
⊢ ∀ (p q : Prop), (p → q) → p → q ∧ p
suggestion: exact fun p q a a_1 => ⟨a a_1, a_1⟩  (closes the goal)
⊢> intro p q h hp
p q : Prop
h : p → q
hp : p
⊢ q ∧ p
suggestion: exact ⟨h hp, hp⟩  (closes the goal)
⊢> exact ⟨h hp, hp⟩
All goals accomplished
finish with :qed [NAME], inspect with :script
⊢> :qed mp_and
saved: theorem mp_and : ∀ p q : Prop, (p → q) → p → q ∧ p
\end{session}

Accepting a suggestion is just typing it.  On a conjunction of two
arithmetic facts, the opening suggestion is already the whole chained
proof, and the session is two lines long:

\begin{session}
λ> :prove (∀ a b : Nat, a + b = b + a) ∧ (∀ a b : Nat, a + 2 * b = b + b + a)
entering prove mode — type tactics; :help for commands
⊢ (∀ (a b : Nat), a + b = b + a) ∧ ∀ (a b : Nat), a + 2 * b = b + b + a
suggestion: constructor <;> omega  (closes the goal)
⊢> constructor <;> omega
All goals accomplished
finish with :qed [NAME], inspect with :script
⊢> :qed two_omegas
saved: theorem two_omegas : (∀ a b : Nat, a + b = b + a) ∧ (∀ a b : Nat, a + 2 * b = b + b + a)
\end{session}

That chain was suggested because \emph{both} halves fall to
\code{omega}.  When they need different tools, the suggestion adapts
to whichever goal is in focus --- arithmetic for the first half, a
found library lemma for the second:

\begin{session}
λ> :prove (∀ a b : Nat, a + b = b + a) ∧ (∀ a b : Nat, a * b = b * a)
entering prove mode — type tactics; :help for commands
⊢ (∀ (a b : Nat), a + b = b + a) ∧ ∀ (a b : Nat), a * b = b * a
suggestion: constructor  (splits into 2 goals)
⊢> constructor
— goal 1 of 2 —
case left
⊢ ∀ (a b : Nat), a + b = b + a
— goal 2 of 2 —
case right
⊢ ∀ (a b : Nat), a * b = b * a
suggestion: omega  (closes the goal)
⊢2> omega
case right
⊢ ∀ (a b : Nat), a * b = b * a
suggestion: exact fun a b => Nat.mul_comm a b  (closes the goal)
⊢> exact fun a b => Nat.mul_comm a b
All goals accomplished
finish with :qed [NAME], inspect with :script
⊢> :qed add_mul_comm
saved: theorem add_mul_comm : (∀ a b : Nat, a + b = b + a) ∧ (∀ a b : Nat, a * b = b * a)
\end{session}

The probes are shaped by the hypotheses as much as by the goal.  An
existential hypothesis draws an \code{obtain} suggestion that names
its witness; once the conjunction inside is the thing in the way, the
suggested step is the \code{obtain} that splits it, chained with the
\code{exact} it enables:

\begin{session}
λ> :prove (∃ x : Nat, 2 < x ∧ x < 4) → ∃ y : Nat, 2 < y
entering prove mode — type tactics; :help for commands
⊢ (∃ x, 2 < x ∧ x < 4) → ∃ y, 2 < y
suggestion: intro h
⊢> intro h
h : ∃ x, 2 < x ∧ x < 4
⊢ ∃ y, 2 < y
suggestion: obtain ⟨x, h1⟩ := h
⊢> obtain ⟨x, h1⟩ := h
x : Nat
h1 : 2 < x ∧ x < 4
⊢ ∃ y, 2 < y
suggestion: obtain ⟨h, h2⟩ := h1 <;> exact Exists.intro x h  (closes the goal)
⊢> obtain ⟨h, h2⟩ := h1 <;> exact Exists.intro x h
All goals accomplished
finish with :qed [NAME], inspect with :script
⊢> :qed ex_two_lt
saved: theorem ex_two_lt : (∃ x : Nat, 2 < x ∧ x < 4) → ∃ y : Nat, 2 < y
\end{session}

A suggestion is not limited to one step.  Here the goal needs
structural induction on a user-defined function; once \code{intro}
brings the variable into the context, the suggested chain arrives as a
complete verified proof:

\begin{session}
λ> def myLen : List Nat → Nat
…>   | [] => 0
…>   | _ :: t => myLen t + 1
…>
λ> :prove ∀ l : List Nat, myLen l = l.length
entering prove mode — type tactics; :help for commands
⊢ ∀ (l : List Nat), myLen l = l.length
suggestion: intro l
⊢> intro l
l : List Nat
⊢ myLen l = l.length
suggestion: induction l <;> simp_all [myLen]  (closes the goal)
⊢> induction l <;> simp_all [myLen]
All goals accomplished
finish with :qed [NAME], inspect with :script
⊢> :qed myLen_length
saved: theorem myLen_length : ∀ l : List Nat, myLen l = l.length
\end{session}

The same pattern scales to arithmetic recursion.  The chain below was
assembled from three of the probes at once: \code{induction} because
the goal mentions a data-typed variable, \code{simp\_all} unfolding
the definition the goal is about, and \code{omega} for the arithmetic
the simplifier leaves behind --- a finished proof of a small theorem
about a function defined moments earlier:

\begin{session}
λ> def double : Nat → Nat
…>   | 0 => 0
…>   | n + 1 => double n + 2
…>
λ> :prove ∀ n : Nat, double n = 2 * n
entering prove mode — type tactics; :help for commands
⊢ ∀ (n : Nat), double n = 2 * n
suggestion: intro n
⊢> intro n
n : Nat
⊢ double n = 2 * n
suggestion: induction n <;> simp_all [double] <;> omega  (closes the goal)
⊢> induction n <;> simp_all [double] <;> omega
All goals accomplished
finish with :qed [NAME], inspect with :script
⊢> :qed double_two_mul
saved: theorem double_two_mul : ∀ n : Nat, double n = 2 * n
\end{session}

A \code{sorry} in an ordinary declaration prints its goal and offers
the same loop via bare \code{:prove}. Prove mode also cooperates with
\code{:synth}; \S\ref{sec:synthprove} shows what that buys.

%==============================================================================
\section{Term synthesis: \code{:synth}}\label{sec:synth}

\code{:synth TYPE} answers the question \emph{``write me a term of
this type''}; read through propositions-as-types, that is
\emph{``prove this.''} Sometimes it answers the stronger question
\emph{``show me that no such term exists.''} It is built on the
vendored Djex library, linked in-process. Djex began by merging two
classic Haskell program synthesizers behind one checked synthesis
foundation: \textbf{Djinn}, a complete and terminating proof
search for intuitionistic propositional logic (Dyckhoff's LJT
calculus) that emits programs, and \textbf{Exference}, a ranked
heuristic search under explicit budgets. It has since moved well
beyond both originals. It fixes a long list of their bugs; the worst
of them let a search that had merely approximated a type report that
type uninhabited, and an exhausted search over an approximated space
now counts for nothing.
It implements much better type-class support, with a unified,
validated constraint contract shared by both engines: Exference
resolves givens, superclasses, and instances, while Djinn checks
contexts and synthesizes dictionary-independent terms. And it adds a
growing level of support for rank-N and impredicative types: goal-side
quantifiers open through polarity-aware plan families, quantified
hypotheses are instantiated at a bounded set of sequent-supplied
types, a separate final tail admits closed forall-free monotypes from
the checked query, and retained loaded schemes keep their own
environment-backed family. Impredicative instantiation is admitted under a
guard that never invents a polytype the query did not supply. Djinn's current
context-free hypothesis rule reaches five leading binders without
raising its fixed search caps; chains with six or more stay opaque.
Its singleton, pairwise, triple, quadruple, and bounded quintuple open/opaque
frontiers cover every choice through eleven independent quantified sites
without enumerating a power set. The quintic categories retain at most 512
plans per orientation; a twelve-site six-open/six-opaque choice is the next
deliberate occurrence gap. That scope is still being
improved, and Leant keeps the boundary to the engines narrow on
purpose, so each improvement arrives by version bump. The engines speak Haskell types;
Leant translates Lean goals into their fragment and their answers back
into Lean.

Three rules run through the design:

\begin{itemize}
\item \textbf{The engine is never trusted.} Every candidate is
  re-elaborated by the Lean backend against the exact goal
  (\code{example :\ (T) :=\ term}) before you see it. A rendering bug
  or an engine bug costs a dropped candidate, never a wrong answer.
\item \textbf{Refusals come with reasons.} A goal outside the
  supported fragment is turned away with a note saying what fell
  outside, not quietly mangled into something answerable.
\item \textbf{Negative verdicts are labeled by strength.}
  ``Provably uninhabited'' appears only when the translation was
  complete and lossless; anything weaker is reported as ``no term
  found within bounds,'' which claims nothing.
\end{itemize}

Transcripts in this section are lightly abridged: trailing candidates
beyond the ones shown, and trailing truncation notes, are elided.
Every line that does appear is a verbatim record of a live session.

\subsection{Higher-order plumbing}\label{sec:plumbing}

The sweet spot is the structural fragment: arrows, products, sums,
$\bot$, $\top$, negation, \code{Iff}, and quantifiers over opaque
variables. These are the ``plumbing'' terms one writes constantly.
Free capital identifiers are auto-bound, so quick queries stay quick:

\begin{session}
λ> :synth ((A → B → C) → (A → B) → A → C)
  it1  fun f g x => f x (g x)
λ> :synth (A × (B ⊕ C) → (A × B) ⊕ (A × C))
  it1  fun ⟨x, y⟩ => match y with | .inl z => .inl ⟨x, z⟩ | .inr w => .inr ⟨x, w⟩
λ> :synth (∀ p q : Prop, (p ↔ q) → ¬p → ¬q)
  it1  fun _ _ ⟨_, f⟩ k x => k (f x)
\end{session}

The first is the S combinator; the second distributes a product over a
sum; the third is contraposition across an \code{Iff}, with the
backward implication projected out by the pattern
\code{⟨\_, f⟩}. Binders are named by role (functions \code{f g h},
values \code{x y z}, negations and continuations \code{k}), a small
thing that makes large candidates much easier to read.

The everyday combinators all come out on the first try, and each
first candidate is the term you would have written yourself ---
\code{flip}, composition, \code{curry}, \code{uncurry}, and product
associativity:

\begin{session}
λ> :synth ((a → b → c) → b → a → c)
  it1  fun f x y => f y x
λ> :synth ((b → c) → (a → b) → a → c)
  it1  fun f g x => f (g x)
λ> :synth ((A × B → C) → A → B → C)
  it1  fun f x y => f ⟨x, y⟩
λ> :synth ((A → B → C) → A × B → C)
  it1  fun f ⟨x, y⟩ => f x y
λ> :synth (((A × B) × C) → A × (B × C))
  it1  fun ⟨⟨x, y⟩, z⟩ => ⟨x, ⟨y, z⟩⟩
\end{session}

Candidates are ranked smallest-first and \emph{bound into the session}
as \code{it1}, \code{it2}, \dots, with bare \code{it} the best one.
They are ordinary definitions; you can evaluate them immediately:

\begin{session}
λ> :synth (a → a → a)
  it1  fun _ x => x
  it2  fun x _ => x
λ> #eval it2 "left" "right"
"left"
\end{session}

Both inhabitants of $a \to a \to a$ that matter, and picking one is
just using its name. For programs, as opposed to proof-irrelevant
propositions, the choice between candidates is the whole point.

An \code{Iff} goal is treated as a pair of implications, so both
directions of a logical identity arrive in one anonymous constructor.
De Morgan and the absorption law:

\begin{session}
λ> :synth (∀ p q : Prop, ¬(p ∨ q) ↔ ¬p ∧ ¬q)
  it1  fun _ _ => ⟨fun k => ⟨fun x => k (.inl x), fun y => k (.inr y)⟩, fun ⟨k1, k2⟩ z => match z with | .inl w => k1 w | .inr x1 => k2 x1⟩
λ> :synth (∀ p q : Prop, p ∧ (p ∨ q) ↔ p)
  it1  fun _ _ => ⟨fun ⟨x, _⟩ => x, fun y => ⟨y, .inl y⟩⟩
\end{session}

Under propositions-as-types, the proof terms \emph{are} plumbing, and
the synthesized spellings make that visible: \code{Iff} symmetry is a
swap, and \code{¬(p ∧ q) ↔ (p → ¬q)} is nothing but currying:

\begin{session}
λ> :synth (∀ p q : Prop, (p ↔ q) → (q ↔ p))
  it1  fun _ _ ⟨f, g⟩ => ⟨g, f⟩
λ> :synth (∀ p q : Prop, ¬(p ∧ q) ↔ (p → ¬q))
  it1  fun _ _ => ⟨fun k x y => k ⟨x, y⟩, fun k1 ⟨z, w⟩ => k1 z w⟩
\end{session}

The textbook curiosities are in range too. No proposition is
equivalent to its own negation, and the synthesized proof is the
classic self-application trick --- the term below applies \code{k} to
itself through \code{f} twice, the propositional shadow of
$\omega\,\omega$:

\begin{session}
λ> :synth (∀ p : Prop, (p ↔ ¬p) → False)
  it1  fun _ ⟨k, f⟩ => k (f (fun x => k x x)) (f (fun y => k y y))
\end{session}

\subsection{Programs you already know}\label{sec:monads}

Some types have one sensible inhabitant, and asking for it by type is
quicker than remembering which library corner it lives in. The bind
of the reader monad, then the bind of the state monad:

\begin{session}
λ> :synth ((S → A) → (A → S → B) → S → B)
  it1  fun f g x => g (f x) x
λ> :synth ((S → A × S) → (A → S → B × S) → S → B × S)
  it1  fun f g x => match f x with | ⟨a, b⟩ => g a b
  it2  fun f g x => match f x with | ⟨a, _⟩ => g a x
  it3  fun f g x => match f x with | ⟨a, b⟩ => (match g a x with | ⟨c, _⟩ => ⟨c, b⟩)
\end{session}

The first candidate threads the state correctly. Now look at the
second: it is type-correct and runs \code{g} on the \emph{initial}
state, throwing away whatever \code{f} did to it. That is the classic
state-threading bug, and the type admits it just as happily. Types
alone cannot tell these apart, which is why all candidates are shown
and each is one keystroke from a test run. Continuation-passing style
works too; this is the bind of the continuation monad:

\begin{session}
λ> :synth ((A → R) → R) → (A → (B → R) → R) → (B → R) → R
  it1  fun f g h => f (fun x => g x h)
\end{session}

With the inductive expansion of \S\ref{sec:inductives} the same game
extends to data. Here is \code{Option.bind}, with the lazy
\code{.none} candidate ranked first and the real one second. When in
doubt, run one:

\begin{session}
λ> :synth (∀ a b : Type, Option a → (a → Option b) → Option b)
  it1  fun _ _ _ _ => .none
  it2  fun _ _ x f => match x with | .none => (.none) | .some y => f y
\end{session}

\subsection{Rank-N and impredicative goals}\label{sec:rankn}

A polymorphic hypothesis is not just cargo: Djex instantiates it at
types the goal itself supplies --- including, under a guard, at
\emph{polymorphic} ones. Here the first candidate applies the
identity hypothesis to the whole goal \code{Q → Q}, an impredicative
instantiation:

\begin{session}
λ> :synth ((∀ p : Prop, p → p) → Q → Q)
  it1  fun f => f _
  it2  fun _ x => x
  it3  fun f x => f _ x
\end{session}

And a Church-encoded pair converts into a real conjunction --- the
quantified hypothesis is instantiated once at \code{p} and once at
\code{q}, fed the matching projection each time:

\begin{session}
λ> :synth (∀ p q : Prop, (∀ r : Prop, (p → q → r) → r) → p ∧ q)
  it1  fun _ _ f => ⟨f _ (fun x _ => x), f _ (fun _ y => y)⟩
\end{session}

(Trailing candidates are elided.) Explicit \code{∀} binders ---
leading, nested, trailing, or
interleaved --- are woven into the candidate's lambda automatically,
and uses of quantified hypotheses get placeholder type arguments
wherever Lean needs them (\code{f \_ x}), so bounded rank-N
candidates verify. Full impredicative inhabitation is undecidable;
beyond the guard the answer is ``no term found within bounds'' and
nothing stronger.

Djinn deliberately keeps three specialization families separate. The
historical local family uses goal variables, opened skolems, premise scopes,
and guarded quantified shapes. A final query-closed tail revisits only
context-free schemes embedded in the requested goal and additionally admits
closed, forall-free subtrees from that checked query. Loaded environment
schemes remain a distinct family with retained source identity and may also use
closed subtrees from loaded signatures.

The separation makes this provider-free goal reachable without treating the
abstract \code{Mono} constant as a type variable:

\begin{session}
λ> axiom QueryClosed.Mono : Type
λ> axiom QueryClosed.Token : Type
λ> axiom QueryClosed.Indexed : Type → Type
λ> :synth ((∀ a : Type, (a → QueryClosed.Token) → a →
…>     QueryClosed.Indexed a) → (QueryClosed.Mono → QueryClosed.Token) →
…>     QueryClosed.Mono → QueryClosed.Indexed QueryClosed.Mono)
  it1  fun f => f _
  it2  fun f g x => f _ (fun _ => g x) x
\end{session}

Both Djinn candidates instantiate \code{f} at
\code{QueryClosed.Mono}; the second supplies a different valid constant
callback. Standalone Exference returns the compact
\code{fun f => f \_}, and combined mode retains both Djinn spellings.
The \code{synth-query-closed-rankn} transcript disables live-library premises,
bounds Exference at 128 steps, and checks all three modes through final Lean
4.31 elaboration. A pure all-engine boundary regression covers the same
projection. The complete synthesis-boundary suite now passes 182 tests.

The new structural and nominal plans come after every established family, so
historical result prefixes do not move. They carry historical local,
loaded-scheme, and checked provider premises for composition, but do not
reclassify loaded declarations as local schemes. The tail independently
retains the five-binder, 16-per-scheme, 64-per-family, and 512-attempt caps.
It is positive-only: exhaustion is \code{NoEvidence}, never a refutation. The
detailed boundary is recorded in the
\href{reports/2026-08-09-query-local-closed-monotype-instantiation.md}{query-local
closed-monotype report}.

Context-free instantiation chains extend to five leading binders. Leant
inserts the inferred type arguments which Djinn leaves implicit, then asks
Lean to verify every rendered candidate against the original type:

\begin{session}
λ> axiom FiveBinder.Five : Type → Type → Type → Type → Type → Type
λ> :synth (∀ A B C D E : Type, (∀ a b c d e : Type, FiveBinder.Five a b c d e) → FiveBinder.Five E D C B A)
  it1  fun _ _ _ _ _ x => x _ _ _ _ _
\end{session}

The focused \code{synth-five-binder-rankn} transcript disables live-library
premises and pins that non-lexical source-order candidate through Lean 4.31
elaboration without truncation. Its first query also discovers an active Lean
instance and retains all five named quantified provider applications. Pure
boundary tests separately retain the five-argument function-elimination shape
and exact provider evidence in Djinn, Exference, and combined mode.

The similarly named live \code{synth-quartic-rankn} regression has a subtler
bridge shape than its Lean surface syntax suggests. After the four outer
\code{FAll} binders for \code{Q}, \code{R}, \code{Z}, and \code{M}, each
vacuous \code{forall A B C D E : Type, Q}-shaped component is serialized as
five ordinary arrows from the shared opaque \code{Type} atom to its codomain,
not as another \code{FAll}. Among the eight result leaves, only the four
polymorphic identity leaves remain quantified. The regression
therefore does not exercise the hypothesis-instantiation
guard. Djex \code{f3dd2495} lets standalone Exference forward those complete
scoped arrow schemes through a structurally materialized nested product. The
quartic submodule revision, \code{c0c1a461}, follows it with tuple-goal provenance:
an independently scheduled structural route may use the eager whole-tree
shortcut once, while fields emitted by the shallow alternative remain
shallow. Equivalent recursive structural histories are suppressed without
losing scoped or environment product reuse at depth. Exference admits the
direct, independently checked candidate at search step 30. The unchanged
4096-step/1024-queue live run still ends with a bounded-tail note recording
36{,}475 queue prunes: search continued after the successful candidate and did
not exhaust its ranked tail.

The live \code{synth-quintic-rankn} regression supplies the missing exact
bridge case. It defines
\code{Wide = forall A B C D E F : Type, A × B × C × D × E × F}, so every binder is
non-vacuous and remains an adjacent \code{FAll}. Leant \code{378f866} coalesces
each uninterrupted \code{FAll} spine to one Djex \code{ForallType} binder list
without crossing an \code{FInst}; the original fragment still supplies every
explicitness slot to the renderer. One \code{Wide} therefore contributes one
Djinn occurrence site rather than six nested sites. With Djex
\code{d728719f}, the first goal selects five non-contiguous opaque
\code{Wide} results beside five opened identities at ten sites. The second
selects five separated open identities beside six opaque \code{Wide} results at
eleven sites, exercising the distinct dual plan. Leant \code{80f123a} records
both direct candidates after Lean 4.31 accepts them; neither live search is
truncated. The sentinel now has six binders, preserving the same occurrence
coverage witness beyond the independent five-binder instantiation cap.

Vacuous local providers can retain the same query-supplied choice even when
their result is an opaque Lean type. Here \code{Demo.Token} becomes a fixed
ambient query rigid, not an unresolved inference variable:

\begin{session}
λ> axiom Demo.Token : Type
λ> :synth ((∀ x : Type, x → x) → ({a : Type 1} → Demo.Token) → Demo.Token)
  it1  fun _ x => @x (∀ (a0_0 : Type _), a0_0 → a0_0)
λ> :set synth-engine exference
synth engine: exference
λ> :synth ((∀ x : Type, x → x) → ({a : Type 1} → Demo.Token) → Demo.Token)
  it1  fun f x => f _ (@x (∀ (a0_0 : Type _), a0_0 → a0_0))
\end{session}

This visible specialization is available only for a fully vacuous,
context-free leading provider chain, and it selects only a closed proper type
already present in the checked query. Fixed ambient rigids may occur in the
provider result; a free flexible variable rejects the route. A local value has
no global binder name, so Leant renders a positional \code{@} application. For
a quantified argument it tries binder domains \code{\_}, \code{Type \_}, and
\code{Prop}, in that order, and backend verification keeps the first spelling
which elaborates. Named globals retain source-directed applications. Each
domain lane retains at most twelve site-and-style spellings, for a hard
36-variant bound only when the local quantified argument is domain-sensitive;
ordinary duplicate spellings collapse to the historical smaller group.

An active Lean instance head can establish a provider-local choice even when
the polytype does not occur in the query. In this session the goal contains
only \code{Gap.Token}; the instance for \code{Gap.C} determines the quantified
argument of the exact provider:

\begin{session}
λ> axiom Gap.Token : Type
λ> class Gap.C (a : Type 1) : Prop where witness : True
λ> instance : Gap.C (∀ x : Type, x → x) := ⟨True.intro⟩
λ> axiom Gap.polyGlobal {a : Type 1} [Gap.C a] : Gap.Token
λ> :set synth-engine djinn
synth engine: djinn
λ> :synth Gap.Token
  it1  Gap.polyGlobal («a» := (∀ (a0_0 : _), a0_0 → a0_0))
λ> :set synth-engine exference
synth engine: exference
λ> :synth Gap.Token
  it1  Gap.polyGlobal («a» := (∀ (a0_0 : _), a0_0 → a0_0))
λ> :set synth-engine both
synth engine: both
λ> :synth Gap.Token
  it1  Gap.polyGlobal («a» := (∀ (a0_0 : _), a0_0 → a0_0))
\end{session}

An active head can also determine a higher-kinded binder which is absent from
the provider body. Leant preserves the argument's bounded
\code{Type~->~\dots~->~Type} kind, so both checked engines keep this
constraint-only application visible:

\begin{session}
λ> namespace Higher
λ> axiom Wrap : Type → Type
λ> class VacuousChoice (F : Type → Type) : Prop where witness : True
λ> instance : VacuousChoice Wrap := ⟨True.intro⟩
λ> axiom VacuousToken : Type
λ> axiom vacuous {F : Type → Type} [VacuousChoice F] : VacuousToken
λ> end Higher
λ> :set synth-engine djinn
synth engine: djinn
λ> :synth Higher.VacuousToken
  it1  Higher.vacuous («F» := Higher.Wrap)
λ> :set synth-engine exference
synth engine: exference
λ> :synth Higher.VacuousToken
  it1  Higher.vacuous («F» := Higher.Wrap)
\end{session}

One exact provider may retain several distinct successful instance-head
vectors. In the same transcript, heads for \code{AlternativeChoice Wrap} and
\code{AlternativeChoice (Pair Nat)} produce exactly
\code{Higher.alternative («F» := Higher.Wrap)} and
\code{Higher.alternative («F» := (@Higher.Pair Nat))}. Standalone Djinn ranks
\code{Wrap} first, standalone Exference ranks \code{Pair Nat} first, and
combined mode uses Djinn's order after stable exact-spelling deduplication.
That is an engine-ranking difference: every mode must retain both exact
alternatives once. The transcript also carries one heterogeneous two-binder
vector whose positional kind arities are one and two, and requires
\code{Higher.multiVacuous («F» := Higher.Wrap) («G» := (@Higher.Triple Nat))}
in all three modes.

Discovery opens no more than five leading type binders on one exact
provider while retaining its erased instance constraints. It inspects at most
32 active heads per provider in Lean resolver order. Each selected head and its
complete constraint closure run under restored metavariable state. The seed
head stays fixed while its returned instance subgoals and every other provider
constraint are solved in that same context; a failure or unresolved binder
rejects the head as a whole. One success contributes one complete vector in
leading-binder order, never an uncorrelated scalar pool. For each argument,
discovery also reduces its type to the admitted first-order kind language and
retains the number of \code{Type} arrow domains. Zero denotes proper type;
positive counts preserve bare and partially applied constructors.

Leant retains at most 16 alpha-distinct vectors per provider and passes no more
than 32 provider/vector associations to Djex. Each vector has exact provider
arity and at most five arguments. The command takes that aggregate prefix
before an argument affects family planning, rigidity, or translation. Live
provider metadata uses
\code{(instantiations (args (kinded N ...) ...))}. Leant folds each bounded
arrow count into a Djex \code{GroundKind}, pairs it with the translated type,
and calls \code{runDjinnQueryWithKindedInstantiationAssignments} or
\code{runExferenceQueryWithKindedInstantiationAssignments}. Live discovery,
wire parsing, and engine filtering all admit at most 64 \code{Type}-arrow
domains, whose simple right-associated kind has \code{2 * 64 + 1 = 129}
constructors. After an assignment passes Djex's provider, scheme, context, and
exact-arity checks, the pinned adapters productively preflight all of that
assignment's supplied kinds before recursive kind inference, same-provider
comparison, kind conversion, or paired type elaboration. Oversized and cyclic
caller-built kinds therefore fail finitely. Metadata-free and
binder-only inventories remain valid. Historical exact arguments stay in their
original vectors and default each position to proper kind; the
\code{(candidates ...)} form parses as proper-kind unary vectors only. A vector
containing depth truncation (\code{FDepth}) or an instance binder
(\code{FInst}) is discarded in full, including an instance binder revealed by
reducing an alias. Provider-prefix fallback keeps metadata attached to its
exact declaration. The checked runners validate provider identity, arity,
supplied positional kinds, closure, and context before consuming each vector
once without Cartesian reconstruction. Thus a later or alpha-identically typed
provider cannot donate its assignment to an earlier one.

Leant also retains applications whose arguments are proper types. A bound
first-order constructor remains a higher-kinded variable in Djex; an opaque
Lean family receives one rigid nominal head across its occurrences. Both
engines can therefore see a query-supplied polytype without learning any hidden
implementation:

\begin{session}
λ> axiom Wrap : Type 1 → Type
λ> :synth ((∀ a : Type 1, Wrap a) → Wrap (∀ b : Type, b → b))
  it1  fun x => x _
λ> :set synth-engine exference
synth engine: exference
λ> :synth (∀ (F : Type 1 → Type), (∀ a : Type 1, F a) → F (∀ b : Type, b → b))
  it1  fun _ x => x _
\end{session}

The serializer admits only first-order kind arrows between universes and only
application arguments whose own type is a universe. Term-indexed families such
as \code{P 3} remain opaque. Unsaturated structural built-ins
\code{And}, \code{Prod}, \code{PProd}, \code{Or}, \code{Sum}, \code{PSum},
\code{Iff}, and \code{Not} remain excluded as higher-kinded assignment heads
until their unsaturated renderer identity is modeled; saturated uses keep
their structural translation. A retained nominal application can support a
positive, Lean-verified candidate but always poisons a negative verdict because
its constant hides structure. Query-derived candidate types come only from
proven proper-type positions below arrows and tuples in the sequent, including
quantified subtrees already present there; provider-local assignment arguments
may also come from the bounded active-instance channel above. This is guarded,
evidence-directed impredicativity, not general second-order search. The
query-derived visible-provider route is capped at five binders and 32 argument
combinations; Djex's internal Haskell-instance route remains monotype-only. The
external Lean evidence route has the separate five-binder, 32-head,
16-vector-per-provider, and 32-vector limits. Open, incomplete,
wrong-arity, or context-bearing vectors fail closed, and chains with six or
more leading binders remain a deliberate incomplete boundary. The original
local-provider contract, scalar active-instance extension, and correlated
follow-up are recorded in the
\href{reports/2026-08-01-scoped-quantified-local-providers.md}{scoped
local-provider report} and the
\href{reports/2026-08-05-provider-local-instance-head-evidence.md}{provider-local
instance-head report}, followed by the
\href{reports/2026-08-05-correlated-instance-head-assignments.md}{correlated
assignment report}, with the shared successor recorded in the
\href{reports/2026-08-09-five-binder-instantiation.md}{five-binder integration
report} and the additive local tail in the
\href{reports/2026-08-09-query-local-closed-monotype-instantiation.md}{query-local
closed-monotype report}.

Unit regressions pin the kinded wire format, kind and argument order, finite
bounds, whole-vector deduplication, and the exact vacuous applications under
Djinn, Exference, and combined mode. The live
\code{synth-provider-higher-kind-assignment} transcript retains the mixed
higher-kinded/rank-N vector, the heterogeneous arity-one/arity-two
multi-vacuous vector, and both exact \code{Wrap} and \code{Pair Nat}
alternatives for one provider. Engine-specific ordering remains ranking policy.

\subsection{Impossibility, proved}\label{sec:negative}

When Djinn's complete search exhausts a fully translated goal, failure
is a theorem:

\begin{session}
λ> :synth (∀ a b : Type, a → b)
provably uninhabited — no closed term of this polymorphic type exists
\end{session}

No Lean tactic gives you this answer (a failing \code{exact?} proves
nothing). Note the careful wording: the verdict is about \emph{closed
terms of the polymorphic type}. It does not say any particular
instance $a \to b$ is empty (choose $a = b$ and it isn't), and
for propositions it is about \emph{constructive} provability, a
distinction \S\ref{sec:classical} exploits. When the translation had
to hide structure behind an opaque atom, the verdict backs off to
match:

\begin{session}
λ> :synth (∀ a : Type, Sigma (fun _ : Nat => a))
no term found within bounds (opaque atoms `Nat` hide structure the engine cannot analyze — not a refutation)
\end{session}

\subsection{Inductive types}\label{sec:inductives}

A non-recursive, non-indexed inductive or structure applied to all its
parameters is expanded into a generalized sum of products:
constructors become introduction rules, case analysis the elimination
rule. This covers the standard library's little workhorses ---

\begin{session}
λ> :synth Ordering
  it1  .lt
  it2  .eq
  it3  .gt
λ> :synth (∀ a : Type, Option (Option a) → Option a)
  it1  fun _ _ => .none
  it2  fun _ x => match x with | .none => (.none) | .some y => y
λ> :synth (∀ e a b : Type, Except e a → (a → b) → Except e b)
  it1  fun _ _ _ x f => match x with | .error y => .error y | .ok z => .ok (f z)
\end{session}

--- \code{Option.join} and \code{Except.map} synthesized rather than
remembered. It extends to \code{Type}-valued classes-as-data like
\code{Decidable}, where the eliminations write themselves, up to and
including the instance combinators a library author would otherwise
write by hand:

\begin{session}
λ> :synth (∀ p : Prop, Decidable p → Decidable (¬ p))
  it1  fun _ x => match x with | .isFalse k => .isTrue k | .isTrue y => .isFalse (fun k1 => k1 y)
λ> :synth (∀ p q : Prop, Decidable p → Decidable q → Decidable (p ∧ q))
  it1  fun _ _ x y => match x with | .isFalse k => .isFalse (fun ⟨z, _⟩ => k z) | .isTrue w => (match y with | .isFalse k1 => .isFalse (fun ⟨_, x1⟩ => k1 x1) | .isTrue x2 => .isTrue ⟨w, x2⟩)
λ> :synth (∀ p q : Prop, Decidable p → Decidable q → Decidable (p → q))
  it1  fun _ _ x y => match x with | .isFalse k => .isTrue (fun z => absurd z k) | .isTrue w => (match y with | .isFalse k1 => .isFalse (fun f => k1 (f w)) | .isTrue x1 => .isTrue (fun _ => x1))
\end{session}

The last two are decidability of conjunction and of implication,
assembled by case analysis on both instance arguments: four
branches each, every branch carrying its own little proof.

Session-declared types participate the moment you declare them, and
refutations over expanded inductives remain sound, because the engine
saw the complete constructor list:

\begin{session}
λ> structure Pair (A B : Type) where
…>   fst : A
…>   snd : B
…>
λ> :synth (∀ a b : Type, a → b → Pair a b)
  it1  fun _ _ x y => ⟨x, y⟩
λ> :synth (∀ a b : Type, Pair a b → Empty)
provably uninhabited — no closed term of this polymorphic type exists
\end{session}

Recursive types cannot be expanded; elimination is where
undecidability lives. Their \emph{constructors} are still sound
introduction rules, though, so building is possible even where
analysis is not:

\begin{session}
λ> :synth Nat
  it1  Nat.zero
  it2  Nat.succ Nat.zero
λ> :synth (∀ a : Type, a → List a)
  it1  fun _ _ => List.nil
  it2  fun _ x => List.cons x List.nil
\end{session}

\subsection{Recursion from the library}\label{sec:synth-library}

\code{:synth} will never invent a \code{Nat.rec}-based program, but
for the everyday recursive types it does not have to: the library
already wrote the recursion. A goal that mentions \code{List} or
\code{Nat} brings a rated inventory of library functions with it ---
\code{List.map}, \code{List.foldr}, \code{List.append},
\code{List.flatten}, \code{List.length}, \code{List.replicate},
\code{Nat.add}, \dots{} --- instantiated at the goal's own types and
offered to the engine as extra premises. The enumeration prefers
proofs that use the goal's own arguments, generally putting a direct
library answer first while retaining distinct choices between
same-typed arguments:

\begin{session}
λ> :synth ((a → b) → List a → List b)
  it1  fun f x => List.map f x
  it2  fun f x => List.reverse (List.map f x)
λ> :synth (List (List a) → List a)
  it1  fun x => List.flatten x
  it2  fun x => List.reverse (List.flatten x)
λ> :synth (Nat → a → List a)
  it1  fun x y => List.replicate x y
  it2  fun x y => List.reverse (List.replicate x y)
λ> :synth (List a → List b → List (a × b))
  it1  fun x y => List.zip x y
  it2  fun _ _ => List.nil
λ> :synth ((a → b → c) → List a → List b → List c)
  it1  fun f x y => List.zipWith f x y
  it2  fun _ _ _ => List.nil
λ> :synth (List a → List a → List a)
  it1  fun x _ => x
  it2  fun _ x => x
  it3  fun x _ => List.reverse x
  it4  fun x y => List.append x y
  it5  fun x y => List.append y x
\end{session}

(Trailing candidates and the truncation notes are elided here.) The
inventory is a ratings list in Djex's \code{*.ratings} format ---
lower is better, 100 or more disables --- and a project file
\code{leant.ratings} (lines of \code{Name Rating}, \code{\#}
comments) merges over the defaults at startup, so re-ranking,
disabling, or growing the inventory is editing a list, not writing
code.

The library search runs beside the constructor search of the previous
section, in a mode where the recursive occurrences are sealed atoms:
without \code{nil}/\code{cons} in play, every candidate must route
through the goal's own arguments and the offered functions, which is
what keeps \code{List.map} ahead of the endless supply of closed
terms that ignore the input. Both searches' survivors are shown
together, library candidates first; \code{:set synth-library off}
restores the constructors-only behavior. As always, the premises are
only \emph{offers}: every candidate is re-elaborated by the backend,
so a useless premise costs search time, never a wrong answer, and
negative verdicts never come from the library run at all.

\subsection{Classical candidates}\label{sec:classical}

Constructive logic goes further than newcomers expect, and
\code{:synth} stays inside it whenever it can. Triple negation
collapses, and contraposition holds in one direction, with no axioms
in sight:

\begin{session}
λ> :synth (∀ p : Prop, ¬¬¬p → ¬p)
  it1  fun _ k x => k (fun k1 => k1 x)
λ> :synth (∀ p q : Prop, (p → q) → ¬q → ¬p)
  it1  fun _ _ f k x => k (f x)
\end{session}

Peirce's law is different: it has no constructive inhabitant, and
\code{:synth} can prove that. But rather than stopping at the
refutation, it first gives constructive live providers a chance to supply a
Lean-verified term. Only if those lanes fail does it ask the classical
question, and answer with a term whose spelling shows what was used:

\begin{session}
λ> :synth (∀ p q : Prop, ((p → q) → p) → p)
  it1  fun _ _ f => match Classical.em _ with | .inl x => x | .inr k => f (fun y => absurd y k)
λ> :synth (∀ p : Prop, ¬¬p → p)
  it1  fun _ k => match Classical.em _ with | .inl x => x | .inr k1 => absurd k1 k
\end{session}

These are the proofs a human would write: one case split on excluded
middle, one \code{absurd} where the contradiction lands. The hard
direction of De Morgan needs excluded middle twice, once per
disjunct, and the candidate reads as exactly that case analysis:

\begin{session}
λ> :synth (∀ p q : Prop, ¬(¬p ∧ ¬q) → p ∨ q)
  it1  fun _ _ k => match Classical.em _ with | .inl x => .inl x | .inr k1 => (match Classical.em _ with | .inl y => .inr y | .inr k2 => absurd ⟨k1, k2⟩ k)
\end{session}

Less famous gaps between the two logics come out the same way. Weak
excluded middle, and the G\"odel--Dummett formula:

\begin{session}
λ> :synth (∀ p : Prop, ¬p ∨ ¬¬p)
  it1  fun _ => match Classical.em _ with | .inl x => .inr (fun k => k x) | .inr k1 => .inl k1
λ> :synth (∀ p q : Prop, (p → q) ∨ (q → p))
  it1  fun _ _ => match Classical.em _ with | .inl _ => (match Classical.em _ with | .inl x => .inl (fun _ => x) | .inr k => .inr (fun y => absurd y k)) | .inr k1 => .inl (fun z => absurd z k1)
\end{session}

Two searches stand behind this. First, one excluded-middle assumption
per atomic subformula is handed to the intuitionistic engine ---
complete for the propositional fragment, since intuitionistic logic
plus atom instances of $p \vee \neg p$ \emph{is} classical
propositional logic. The findings are rendered as \code{Classical.em}
case splits. If that misses, the goal's double-negation translation runs
instead (complete by Glivenko's theorem) and the candidate is wrapped
in \code{Classical.byContradiction}. Both kinds are verified like any
other candidate. The whole behavior sits behind a switch:

\begin{session}
λ> :set synth-classical off
synth classical: off
λ> :synth (∀ p : Prop, p ∨ ¬ p)
provably uninhabited — no closed term of this polymorphic type exists
(constructively — a classical proof may still exist; this is not a disproof of the proposition)
\end{session}

\subsection{Dependent formulas as cargo}\label{sec:dependent}

Leant's translation makes no attempt to analyze dependent types.
It does not refuse them either: a dependent subformula becomes an
opaque \emph{atom}, compared up to $\alpha$-equivalence, which the
engine can carry from one side of a goal to the other. Transport
works; analysis does not:

\begin{session}
λ> opaque P : Nat → Prop
λ> opaque Q : Prop
λ> :synth ((∀ n : Nat, P n) ∧ Q → Q ∧ (∀ n : Nat, P n))
  it1  fun ⟨x, y⟩ => ⟨y, x⟩
\end{session}

The engine never looked inside $\forall n,\, P\,n$; it swapped a
sealed box. A goal that would require opening the box (an induction, a
rewrite, a case split on an index) is refused with a reason; that work
belongs to \code{:prove}.

\subsection{Synthesis inside a proof}\label{sec:synthprove}

Bare \code{:synth} in prove mode targets the current goal \emph{with
its hypotheses as premises}. Candidates are offered as \code{it1},
\code{it2}, \dots\ exactly as in the session, except that here
\code{itN} splices the candidate applied to your hypotheses, so
\code{exact it1} closes the goal. The hypotheses can carry real
weight; below, the case analysis on the \code{Decidable} instance
comes from \code{hd}:

\begin{session}
λ> :prove ∀ p q : Prop, Decidable p → (p → q) → (¬p → q) → q
entering prove mode — type tactics; :help for commands
⊢ ∀ (p q : Prop) (a : Decidable p), (p → q) → (¬p → q) → q
suggestion: exact fun p q a a_1 a_2 => dite p a_1 a_2  (closes the goal)
⊢> intro p q hd hpq hnq
p q : Prop
hd : Decidable p
hpq : p → q
hnq : ¬p → q
⊢ q
suggestion: exact dite p hpq hnq  (closes the goal)
⊢> :synth
(synthesizing with hypotheses p q hd hpq hnq as premises)
  it1  fun _ _ x f g => match x with | .isFalse k => g k | .isTrue y => f y
⊢> exact it1
All goals accomplished
finish with :qed [NAME], inspect with :script
⊢> :qed byCases
saved: theorem byCases : ∀ p q : Prop, Decidable p → (p → q) → (¬p → q) → q
\end{session}

Note the contrast with the tactic suggestion above it: the suggestion
\emph{found} the library's \code{dite}, while \code{:synth}
\emph{composed} the case split from the goal's own material. Both are
verified; they simply answer different questions.

It also slots into an ordinary multi-goal proof. Split a conjunction
with \code{constructor} and let \code{:synth} deal with each half in
turn. Notice the prompt tracking the number of open goals, and the
projections from \code{h} showing up without being asked for:

\begin{session}
λ> :prove ∀ p q r : Prop, (p → q) ∧ (q → r) → p → r ∧ (p → q)
entering prove mode — type tactics; :help for commands
⊢ ∀ (p q r : Prop), (p → q) ∧ (q → r) → p → r ∧ (p → q)
suggestion: intro p q r h h1 <;> simp_all  (closes the goal)
⊢> intro p q r h hp
p q r : Prop
h : (p → q) ∧ (q → r)
hp : p
⊢ r ∧ (p → q)
suggestion: simp_all only [forall_const, imp_self, and_self]  (closes the goal)
⊢> constructor
— goal 1 of 2 —
case left
p q r : Prop
h : (p → q) ∧ (q → r)
hp : p
⊢ r
— goal 2 of 2 —
case right
p q r : Prop
h : (p → q) ∧ (q → r)
hp : p
⊢ p → q
suggestion: simp_all only [forall_const]  (closes the goal)
⊢2> :synth
(synthesizing with hypotheses p q r h hp as premises)
  it1  fun _ _ _ ⟨f, g⟩ x => g (f x)
⊢2> exact it1
case right
p q r : Prop
h : (p → q) ∧ (q → r)
hp : p
⊢ p → q
suggestion: simp_all only [forall_const, imp_self]  (closes the goal)
⊢> :synth
(synthesizing with hypotheses p q r h hp as premises)
  it1  fun _ _ _ ⟨f, _⟩ x _ => f x
⊢> exact it1
All goals accomplished
finish with :qed [NAME], inspect with :script
⊢> :qed chain
saved: theorem chain : ∀ p q r : Prop, (p → q) ∧ (q → r) → p → r ∧ (p → q)
\end{session}

Unlike \code{exact?}, which finds an \emph{existing} lemma, this
composes a new term from the goal's own material: a constructive
complement to the finisher tactics, needing no premise database and
no imports. The division of labor that
emerges is pleasant: the genuinely dependent moves (quantifiers,
induction, rewriting) stay yours, and the propositional bookkeeping
between them does not.

The classical fallback of \S\ref{sec:classical} follows \code{:synth}
into prove mode. Double-negation elimination has no constructive
proof, so the candidate arrives as an excluded-middle case split, and
\code{:qed} turns it into a named theorem --- no tactic beyond
\code{exact} was typed:

\begin{session}
λ> :prove ∀ p : Prop, ¬¬p → p
entering prove mode — type tactics; :help for commands
⊢ ∀ (p : Prop), ¬¬p → p
suggestion: exact fun p a => Classical.byContradiction a  (closes the goal)
⊢> :synth
  it1  fun _ k => match Classical.em _ with | .inl x => x | .inr k1 => absurd k1 k
⊢> exact it1
All goals accomplished
finish with :qed [NAME], inspect with :script
⊢> :qed not_not_elim
saved: theorem not_not_elim : ∀ p : Prop, ¬¬p → p
\end{session}

\subsection{Engines and budgets}\label{sec:engines}

The default engine is Djinn: complete, terminating, and the only
source of refutations. \code{:set synth-engine exference} switches to
the ranked heuristic search, and \code{both} runs the two together ---
negative verdicts still come only from Djinn. Standalone lanes send at
most 12 fresh candidate groups to Lean. A combined lane gets 24 and
retains both standalone frontiers: writing D and E for fresh Djinn and
Exference groups, its order is \code{D1–D4, E1–E12, D5–D12}, then
alternating tails. Within each Exference invocation, Leant
stable-deduplicates rendered groups before applying the internal 60-candidate
collection window. The first spelling remains authoritative, while repeated
backend histories cannot hide later distinct terms behind the bound; the
outer 12- and 24-group verification frontiers then apply. Exference runs under
explicit budgets
(\code{:set synth-steps}, default 4096), and truncation is always
reported:

\begin{session}
λ> :set synth-engine both
synth engine: both
λ> :synth ((A → B → C) → (A → B) → A → C)
  it1  fun f g x => f x (g x)
note: exference: search truncated: step limit reached, queue limit pruned 31990
\end{session}

Every engine mode first gives a structurally accepted goal a
provider-free lane and asks Lean to verify its candidates. When no
provider-free term verifies, Lean discovers at most 80 live term providers
under the namespaces used by the target (plus exact session declarations).
A complete Djinn refutation is retained as a sound fallback while those
constructive provider lanes run: the first candidate accepted by Lean wins;
an empty, unavailable, timed-out, or unsuccessful provider search restores
the refutation before any classical retry.
The focused \code{synth-provider-refutation-fallback} transcript checks exact
rank-N overrides under Djinn, Exference, and combined mode. It also enables
classical search for a Peirce-shaped goal and confirms that an exact
constructive provider wins first, while a no-provider control retains the
original sound refutation. Pure tests pin the same direct, rank-N, and
unusable-provider boundaries below the REPL layer.
Type constructors and type families, whose fully peeled result is a
sort, are excluded before search. Exact-result session and public
declarations lead the ranking;
conventional implementation-worker names remain eligible but follow
public fallbacks, while session declarations bypass that naming
heuristic. Exference searches the ranked inventory directly; Djinn
tries discovery-order prefixes of 1, 4, and 16 providers before the full
bounded inventory, and may specialize a loaded scheme at a closed
monotype or guarded rank-N polytype. The serializer derives a canonical
query from the
sorted, deduplicated namespace roots and final result head, rather than
from the goal's raw spelling. Successful inventories --- including an
empty one --- are kept in a generation-aware 12-entry LRU. Imports,
ordinary declarations, and session-restoring operations invalidate that
generation; generated \code{it1}, \code{it2}, \dots\ bindings cannot be
providers and preserve it.

For an exact provider with erased instance constraints, the same discovery
pass can attach bounded active-instance-head assignments. Resolution follows
Lean's own instance order and isolates the selected head plus its complete
constraint closure. One successful head yields one ordered vector; an
incomplete or unusable head yields nothing. Provider slicing retains the
vector only with its source declaration; after whole-vector context/depth and
arity filtering Leant takes the command-wide prefix before planning or
translation. Each argument crosses with its bounded \code{Type}-arrow kind;
Leant reconstructs a Djex \code{GroundKind} and calls the checked kinded Djinn
or Exference assignment runner. Leant rejects residual kind arities above 64
before that bridge, and pinned Djex independently rejects a supplied
\code{GroundKind} above 129 constructor nodes before recursive operations on
that assignment.
Lean still reconstructs every dictionary
during final elaboration---only the closed kind/type vector crosses the bridge.

Exact rendered spellings keep their first scheduled occurrence. Before a
later provider lane is forced or capped, spellings already rejected by Lean
are removed from each engine's source stream and newly empty groups are
dropped. Rediscovered failures therefore spend none of that lane's 12- or
24-group fresh frontier.

The pure searches answer in microseconds to milliseconds; the cost
center is backend verification, a few hundred milliseconds per
candidate. A wall-clock guard (default 20\,s,
\code{LEANT\_SYNTH\_TIMEOUT}) covers the quantified goals whose
bounded instantiation can widen the space; hitting it is reported as
``no answer,'' never as a verdict.

\subsection{What \code{:synth} will not do}\label{sec:limits}

Some plain statements of present limits. It will not synthesize recursion or
induction. Complete compatible recursive families contribute bounded positive
construction to Djinn and one constructor layer of elimination to Exference;
recursive fields are never opened again. It will not analyze dependent
structure; atoms are transported, never opened. Opaque and bound-variable
proper-type applications retain their ordered arguments, and qualifying
inductives share one exact parametric family across distinct
\code{Option}/\code{List} occurrences. Term-indexed, dependent,
ambiguous, or incompatible families keep their conservative fallback.
Quantifier handling beyond the propositional fragment follows
Djex's bounded rank-N and guarded impredicativity rules, which are
sound and terminating over an undecidable space and widen as Djex
improves. The current plan family is exhaustive through eleven
independent quantified sites; a twelve-site six-open/six-opaque goal
is the next deliberate occurrence-plan gap, while hypothesis
instantiation chains with six or more binders remain opaque. Such
higher-rank goals answer ``no term
found within bounds'' and nothing stronger. And although Djex can resolve
type-class evidence itself, Leant leaves instance arguments to Lean's
elaborator during verification, which sits nearer the instance tables
anyway. A non-dependent instance-implicit goal binder is retained as a
render-only slot: it contributes no dictionary premise to either engine, but
Leant inserts its wildcard before later lambda arguments and withholds complete
negative evidence across the erased boundary. Constrained rank-N hypotheses
therefore keep their instance applications implicit without shifting the
candidate's term binders. Provider discovery now consults active Lean instance
heads only to obtain bounded ground-kinded assignments for the exact provider
whose erased constraints they collectively solve. Constraint-only vacuous
higher-kinded binders are supported when each argument kind is a bounded
\code{Type~->~\dots~->~Type} chain; unsaturated structural built-in heads
remain excluded. Leant never serializes
dictionaries or general instance resolution into Djex. A selected head is
accepted only after its own subgoals and every remaining provider constraint
close in one isolated metavariable context. Up to 32 heads are inspected and
16 distinct complete vectors retained per provider, with at most 32 vectors
passed in total; open, incomplete, wrong-arity, depth-truncated, and contextual
vectors fail closed. The reducible
control \code{Gap.Contextual := \{a : Type\} → [Inhabited a] → a} is therefore
excluded even when an active class instance names it. This provider-local
channel can expose a polytype absent from the query, but it is bounded checked
evidence rather than general impredicative inference. The fragment is bounded
on purpose: inside it, answers are fast and verified. Refutations are complete
for the provider-free structural calculus but provisional with respect to the
live environment. Leant searches the bounded, best-effort provider inventory
before reporting one: a verified provider candidate overrides it, while an
empty or unsuccessful provider search preserves it as the original sound
fallback. The resulting verdict is not an exhaustive theorem about every
declaration or axiom in the Lean environment.

%==============================================================================
\section{How it works, briefly}\label{sec:how}

A \code{:synth} query passes through five checked stages. A Lean
metaprogram (compiled once into a cached side environment) elaborates
the goal and serializes it into the engines' fragment: structural
connectives are decomposed; qualifying inductives expand into
constructor lists; proper-type applications retain rigid or bound heads
and ordered arguments; everything else becomes an $\alpha$-normalized
opaque atom. The fragment translator either accepts the result or
refuses with a reason. The cached synthesis environment remembers the
session history it has replayed: an unchanged history is reused, an
append replays only its new suffix, and undo or another non-prefix
change rebuilds from the cached import-and-serializer base. Generated
result bindings participate in this replay so later goals can mention
them without flushing the provider cache. The selected engine searches;
its candidates are rendered back into Lean syntax, with constructor
names restored, binders named by role, and anonymous-constructor and
leading-dot spellings preferred over fully qualified fallbacks. Finally the
backend re-elaborates each candidate against the original goal, and
only survivors are shown and bound. The engine boundary is narrow by
design: the searches behind it are swappable, and nothing they emit is
believed until Lean has type-checked it.

An unpickled environment is an explicit history barrier. Leant copies it to
a private process-lifetime artifact so backend restart can restore that base
before replaying the commands entered afterward; the reconstructed undo stack
therefore still bottoms out at the snapshot. \code{:reset} and \code{:load}
leave snapshot mode, whereas \code{:import} requires that explicit transition
because imports cannot be added to an opaque Lean environment. A Leant
\code{:pickle} also writes a content-fingerprinted, serializer-ABI-keyed
synthesis companion. Foreign or stale sidecars are ignored without rejecting
the main upstream snapshot; if that snapshot exposes Lean's metaprogramming
API, Leant compiles current synthesis tooling directly over it.

%==============================================================================
\section{Status}\label{sec:status}

Leant is a working prototype under active development.
The synthesis design and its phased history are documented in
\code{docs/SYNTHESIS\_PROPOSAL.md}; transcript-level golden tests live in
\code{test/}. The project began inside the \emph{ProveIt} repository
and was split out with its history intact. Expect sharp edges: the project exists to find out how far a
conversational, verification-gated synthesis workflow can be pushed,
and it moves at the speed of that question.
\end{document}
